\section{Challenges and Future}
\label{sec:challenges}

\subsection{Robust Diagnostic Evaluation}
\label{sec:eval-diagnostic}

\citet{chen2024x2dfd} and \citet{sun2025edvd} argue that classification-centric metrics such as clip-level \textit{AUC}/\textit{EER} are often not sufficiently evidential or interpretable for security-sensitive settings.
Under our \emph{factual-fidelity} objective in \S\ref{sec:scope}, strong closed-set discrimination does not necessarily mean the detector can substantiate \emph{which} real-world proposition is violated, \emph{where} the violation occurs, or \emph{why} the judgment is reliable.
Evaluation should therefore be \emph{diagnostic}: it should probe the detector's reliance on complementary pathways (perceptual/temporal cues vs.\ claim-level semantic verification), and reveal the failure modes that lead to factual-fidelity breakdown.

Concretely, one critical axis is robustness of perceptual and temporal evidence under distribution shift.
As diffusion and large video models increasingly suppress noticeable per-frame artifacts, evaluation must stress-test video-level consistency and robustness-to-shift rather than only in-domain performance.
This pressure is already visible in explainable detection settings such as EDVD-LLaMA~\citep{sun2025edvd} and DeeptraceReward~\citep{Fu2025DeeptraceReward}, in generator-diverse detection benchmarks such as GenVidBench~\citep{ni2025genvidbench}, and even in physics-oriented generative evaluation such as Physics-IQ~\citep{motamed2025generative}.
\citet{DBLP:journals/corr/abs-2601-05986} and \citet{hasan2026uneven} further show that strong clean-set scores can mask severe brittleness under transfer-based adversarial attacks and modern synthesis pipelines, leaving a wide gap between benchmark discrimination and deployment robustness.
This motivates cross-dataset generalization protocols, systematic perturbation sweeps over codec, bitrate, and resolution, and deployment-oriented reporting at fixed operating points such as \textit{TPR@FPR=$\alpha$} to better reflect real risk.

A second axis is whether the system can treat a video as making checkable claims and verify them with grounded evidence.
Beyond synchronization, cross-modal alignment tests for retrieval and temporal localization should be complemented with fact-level and knowledge-level verification.
Common sense reasoning for deepfake detection~\citep{zhang2024common} pushes evaluation toward explicit question answering, LVLM-DFD~\citep{Yu2025LVLMDFD} and AuthGuard~\citep{Shen2025AuthGuard} emphasize grounded multimodal reasoning, and BusterX++~\citep{Wen2025BusterXpp} adds unified cross-modal detection-and-explanation settings. These developments also motivate measuring rationale quality and grounding, not just label accuracy.
Together, these diagnostics align evaluation with the shift from perceptual-fidelity discrimination to factual-fidelity verification.

\subsection{Claim-Level and Dynamic Evaluation}
\label{sec:eval-evidence}

A natural next step is to make benchmarks \emph{evidence-first} by moving from clip-level labels to claim-level supervision.
In this setting, each clip is decomposed into a compact set of propositions specifying who, when, where, and what happens, and each proposition is paired with timestamped in-video evidence.
Evaluation then scores both claim correctness and evidence grounding, making ``what is fake'' and ``where it is fake'' traceable rather than implicit.
This direction is consistent with recent explainable AIGC detection resources that augment binary supervision with rationales, traces, and defect-level annotations.
X2-DFD~\citep{chen2024x2dfd} and EDVD-LLaMA~\citep{sun2025edvd} make explanation part of the detector output, Ivy-Fake~\citep{jiang2025ivy} and DAVID-XR1~\citep{gao2025david} attach rationales and defect-level traces to AI-generated video analysis, DeeptraceReward~\citep{Fu2025DeeptraceReward} adds human-perceived fake traces, and ExDDV~\citep{hondru2025exddv} turns explainable supervision into a benchmark design target.

To prevent shortcut learning and isolate genuine reasoning, evidence-first benchmarks should also incorporate targeted stress tests.
One complementary approach is claim-level stress testing for event logic: construct counterfactual or adversarial cases by rewriting event scripts or state transitions while controlling visual-quality confounds, so performance differences primarily reflect relational and event reasoning rather than superficial appearance cues. This concern is explicit in Can We Leave Deepfake Data Behind in Training Deepfake Detector?~\citep{cheng2024can} and FakeChain~\citep{heo2025fakechain}, both of which expose how easily detectors can rely on shallow cues.
In addition, benchmarks increasingly include supervision for tampered locations, spatiotemporal traces, and grounded explanations. EDVD-LLaMA~\citep{sun2025edvd}, DeeptraceReward~\citep{Fu2025DeeptraceReward}, ExDDV~\citep{hondru2025exddv}, and DAVID-XR1~\citep{gao2025david} all move beyond binary labels in this direction, while Video Reality Test~\citep{wang2025videorealitytestaigenerated} adds realistic high-fidelity content where obvious artifacts are weak.
Where available, cross-validating content-side decisions with provenance and authentication signals can further improve trustworthiness in security-sensitive settings.

Finally, given fast-evolving generators, we expect evaluation to move beyond static test sets toward a closed loop between a \emph{real-time detection platform} and a \emph{dynamic benchmark}.
In an arena- or leaderboard-style regime, the benchmark is continuously refreshed with outputs from newly released generators and with realistic ingestion or transcoding effects such as codec and resolution changes; detectors are then periodically re-evaluated under a unified protocol to produce regression trajectories rather than one-off scores~\citep{wang2025videorealitytestaigenerated}. Such a setup measures not only snapshot accuracy, but also robustness decay, adaptation lag, and failure modes under sustained distribution shift. Hard cases that repeatedly fool strong detectors can then be promoted into subsequent benchmark refreshes with failure-type annotations, turning deployment feedback into iterative stress testing rather than passive reporting.


\subsection{Unified Explainable Detection}
\label{sec:dualview-factual}
\citet{wang2025videorealitytestaigenerated} indicates that, for AIGC-V with sparse factual content but high perceptual fidelity and well-aligned modalities, VLMs can perform near random in perceptual-fidelity discrimination, far below humans. This highlights a core gap of language-only detection and the need for visual-view evidence. At the same time, stronger generators and anti-detection pipelines make it increasingly unsafe to assume that authenticity can be decided from perceptual traces alone. Common sense reasoning~\citep{zhang2024common}, generalizable LVLM reasoning~\citep{Yu2025LVLMDFD}, language guidance~\citep{Shen2025AuthGuard}, explainable video reasoning~\citep{gao2025david}, and memory-anchored multimodal reasoning~\citep{chen2025memory} all point toward detectors that must reason over semantic and factual content as well. When perceptual fidelity is high and modalities align, decisive evidence may still come from subtle visual statistics or long-range temporal inconsistencies; frequency-aware blending cues~\citep{zhou2024freqblender}, local dynamic inconsistency learning~\citep{gu2022delving}, natural consistency representations~\citep{zhang2024learning}, and plug-and-play video-level blending with spatiotemporal adapter tuning~\citep{yan2025generalizing} all support this view. Conversely, when generators suppress artifacts and preserve visual coherence, visual-only cues can be insufficient and fact-level reasoning becomes necessary; this is exactly the trajectory emphasized by common-sense and reasoning-heavy detection methods~\citep{zhang2024common,Yu2025LVLMDFD,Shen2025AuthGuard,gao2025david}. Unified detection is therefore a two-pathway problem: perceptual evidence plus fact-level verification.

In practice, this motivates cross-layer evidence fusion: frame-level intrinsic cues from Layer~1, sequence-level motion and physics coherence from Layer~2, multimodal verification from Layer~3, and external-knowledge reasoning from Layer~4 should be integrated into a unified evidence graph, where corroborating evidence boosts confidence and conflicts trigger abstention or human review. The goal, however, is not to average heterogeneous scores, but to preserve evidential granularity so that low-level traces, localized cross-modal conflicts, and claim-level checks remain individually inspectable. The key challenge is to connect low-level cues to high-level conclusions in an auditable way, as X2-DFD~\citep{chen2024x2dfd} and EDVD-LLaMA~\citep{sun2025edvd} already make clear. A practical system should output structured evidence objects, such as suspicious segments and defect-level traces~\citep{gao2025david,anshul2025next}, localized cross-modal mismatches~\citep{katamneni2024contextualcrossmodalattentionaudiovisual,chen2024x2dfd}, and tested claims or entities~\citep{zhang2024common}, rather than only free-form rationales. Tool-augmented pipelines such as LAVID~\citep{Liu2025LAVID} help only when each tool invocation is tied to a concrete sub-claim and yields reusable evidence objects; memory-anchored multimodal reasoning~\citep{chen2025memory} points to the same requirement.

\subsection{Evidence-First Trustworthy Detection}
\label{sec:evidence-first}
To match trustworthy detection that truly targets factual fidelity, system design should follow an Evidence-First principle with an explicit reasoning path: Identify, Localize, and Explain. X2-DFD~\citep{chen2024x2dfd}, DeeptraceReward~\citep{Fu2025DeeptraceReward}, and EDVD-LLaMA~\citep{sun2025edvd} already move in this direction. The key implication is that explanation should be downstream of evidence extraction, not a post-hoc verbal gloss over a classifier score. This path should separate candidate-issue discovery, where-and-when localization, and explanation of which perceptual, temporal, cross-modal, or fact-level constraints are violated. Such a design requires consistent evidence schemas, as DAVID-XR1~\citep{gao2025david} and X2-DFD~\citep{chen2024x2dfd} suggest, together with calibrated uncertainty emphasized by reliability-centered analyses~\citep{Wang2024DeepfakeReliability}. Calibration matters because trustworthy deployment depends on knowing when evidence is weak, conflicting, or incomplete rather than forcing a binary answer.

Within this dual-pathway setup, each tool invocation or knowledge citation should be bound to a specific argumentation step~\citep{Liu2025LAVID}. Content-side analysis should also be cross-validated against source-side provenance and authentication signals when available, rather than treated as an unrelated defense channel. This connection is explicit in media-forensics and defense surveys~\citep{Verdoliva2020MediaForensicsDeepFakes,Deng2025AIGCDefensesSurvey}, in technical provenance standards such as C2PA~\citep{c2pa2024spec}, and in complementary verification frameworks built on face geometry~\citep{tursman2020towards} or blockchain-backed authenticity pipelines~\citep{hajjej2026framework}. This matters beyond standalone detection accuracy: \citet{sagar2026fact} suggests that detector outputs can become misleading priors when they are injected into claim verification pipelines without explicit evidence grounding. The practical difficulty is to reconcile potentially conflicting evidence and surface it in a unified evidence space that supports auditing and abstention under uncertainty~\citep{Wang2024DeepfakeReliability}. A trustworthy system should therefore preserve both agreement and disagreement across evidence sources, rather than collapsing them too early into a single confidence score.

\begin{surveytakeaways}[What Trustworthy Detection Requires]
  \item \textbf{Diagnose the real failure.} Evaluation should reveal where factual fidelity breaks, not only whether real and fake clips separate on a closed set.
  \item \textbf{Supervise evidence, not only labels.} Claim-level annotations, timestamped evidence, shortcut controls, and dynamic testbeds are needed as generators evolve.
  \item \textbf{Return structured evidence, not only scores.} Cross-layer fusion is useful only when cues become traceable claims, segments, entities, or mismatches rather than free-form confidences or rationales.
  \item \textbf{Treat deployment as evidence governance.} Provenance checks, calibration, and abstention are part of trustworthy detection when evidence conflicts or remains incomplete.
\end{surveytakeaways}
